% firlatma_karsilastirma.tex — sahte EDM ölçümünde fırlatma yöntemi: makaledeki yöntem ile donmuş spin yöntemi
% Derleme: xelatex firlatma_karsilastirma.tex (iki kez)
\documentclass[11pt,a4paper]{article}
\usepackage{fontspec}
\usepackage{polyglossia}
\setdefaultlanguage{turkish}
\setotherlanguage{english}
\usepackage{amsmath,amssymb}
\usepackage{graphicx}
\usepackage{booktabs}
\usepackage{array}
\usepackage{tabularx}
\usepackage[margin=2.3cm]{geometry}
\usepackage{xcolor}
\usepackage{xurl}
\usepackage{hyperref}
\usepackage{caption}
\usepackage{float}
\usepackage{enumitem}
\usepackage{fancyvrb}
\usepackage{fvextra}
\setmonofont{DejaVu Sans Mono}[Scale=0.82]
\hypersetup{colorlinks=true,linkcolor=blue!50!black,urlcolor=blue!50!black}
\captionsetup{font=small,labelfont=bf}
\setlength{\parskip}{0.45em}
\setlength{\parindent}{0pt}
\graphicspath{{sonuclar/}}
\newcommand{\kod}[1]{{\XeTeXinterchartokenstate=0\texttt{#1}}}
\newcommand{\BEL}{\textsc{[belgelenmiş]}}
\newcommand{\YOR}{\textsc{[yorum]}}
\newcommand{\HIP}{\textsc{[hipotez]}}
\newcommand{\sekil}[3][0.95]{\begin{figure}[H]\centering\includegraphics[width=#1\textwidth]{#2}\caption{#3}\end{figure}}
\DefineVerbatimEnvironment{kodblok}{Verbatim}{fontsize=\small,frame=leftline,framerule=1.2pt,rulecolor=\color{black!30},xleftmargin=4mm,numbers=none,breaklines=true,breakanywhere=true}
\newcommand{\BaslangicTablo}{%
\begin{center}\small
\begin{tabular}{lcccc}
\toprule
fırlatma hücresi & 0 & 6 & 12 & 18 \\
\midrule
başlangıç noktası, eski kural ($\delta=0$): $f_0$ & 1088 & 2729 & $-456$ & 95 \\
başlangıç noktası: $\omega_0$ [rad/s] & $+1.81$ & $-2.48$ & $-0.39$ & $-0.015$ \\
ikinci nokta: $\omega_1$ [rad/s] & $-0.90$ & $+1.2$ & $+0.19$ & $+0.0075$ \\
ara değerleme: $f^*$ & 82.27 & 79.55 & 82.00 & 81.96 \\
\textbf{$\delta^*$'da doğrudan ölçüm: $f$} & \textbf{82.04} & \textbf{81.87} & -- & -- \\
\quad ve orada $\omega_h$ [rad/s] & $-7\times10^{-6}$ & $+2\times10^{-5}$ & -- & -- \\
enerji kuralından başlayarak (Test 1): $f^*$ & 82.04 & 81.90 & 81.89 & 81.94 \\
\bottomrule
\end{tabular}
\end{center}}
\newcommand{\BaslangicYorum}{%
Okunuşu~\BEL:
\begin{enumerate}[nosep]
\item \emph{Donmuş nokta başlangıç kuralından bağımsız bulunuyor.} Eski kuraldan başlanınca bulunan
$\delta^*$ değerleri ($+2.313\times10^{-7}$, $-3.168\times10^{-7}$, $-5.02\times10^{-8}$, $-1.95\times10^{-9}$),
donmuş spin yörüngesinin o hücredeki $|p|$ sapmasıyla (Test~2'nin çapraz kontrolü:
$+2.313\times10^{-7}$, $-3.168\times10^{-7}$, $-5.024\times10^{-8}$, $-1.952\times10^{-9}$) dört basamak aynı.
Yani iki farklı başlangıç kuralı aynı yörüngeyi buluyor.
\item \emph{$\delta^*$'da doğrudan ölçülen $f$, enerji kuralıyla bulunan değerle \%0.04 içinde aynı}
(82.04 ve 81.87'ye karşı 82.04 ve 81.90); orada $\omega_h$ gerçekten sıfır ($10^{-5}$ rad/s).
\item \emph{İki noktalı ara değerleme ise bu testte \%3'e kadar sapıyor} (hücre 6: 79.55). Nedeni: eski
kuralda $\omega_h$'nin $\delta$'ya eğimi farklıdır ($\approx-7.8\times10^{6}$ rad/s; aracın varsaydığı
$-5.24\times10^{6}$ enerji kuralına aittir). Bu yüzden ikinci nokta donmuş noktaya düşmez
($|\omega_1|\approx1$~rad/s) ve ara değerleme geniş bir aralıkta yapılır. Donmuş noktanın dışında $f$'nin
davranışı tam doğrusal olmadığından ara değerleme hata yapar~\HIP.
\item \emph{Önceki bütün üretim koşumları için sonuç~\BEL:} enerji kuralıyla başlayan 666 koşumun hepsinde
ikinci nokta donmuş noktaya çok yakındı ($|\omega_1|$ en fazla $2.7\times10^{-4}$, medyanı $10^{-5}$ rad/s).
Bu testteki 0.9--1.2 rad/s ile karşılaştırınca, önceki sonuçlar fiilen donmuş noktada doğrudan ölçülmüş
sayılır; ara değerleme hatası ihmal edilebilir.
\item \emph{Araçta yapılan iyileştirme:} $|\omega_1|>10^{-3}$ rad/s ise araç artık $\delta^*$'da üçüncü bir ölçüm
yapıyor ve o değeri raporluyor (\kod{W1\_ESIK}). Önceki koşumların hiçbirinde bu eşik aşılmadığı için eski
sonuçlar değişmez; araç başlangıç kuralına karşı sağlam hâle gelir.
\end{enumerate}}

\newcommand{\HamEskiCW}{$-0.08$}
\newcommand{\HamEskiCCW}{$-313$}
\newcommand{\HamEskiC}{150}
\newcommand{\FarkYorum}{%
Okunuşu~\BEL: fark, kapalı yörüngenin büyüklüğüyle birlikte yaklaşık doğrusal büyüyor (kesikli çizgi
eğim 1). Bu, Test~5'teki açıklamayla tutarlı: hata $\propto\sigma\cdot\sigma$, kapalı yörünge
büyüklüğü (rms$_x\cdot$rms$_y$) de $\propto\sigma^2$. Üç rejim:
\begin{itemize}[nosep]
\item \emph{İki demet düzeltmesi} (yörünge $\sim0.01~\mu$m$^2$): fark $\le0.02$, yani sayısal taban; iki
yöntem uyuşuyor.
\item \emph{Tek demet, düzeltilen CCW yönü} ($\sim1~\mu$m$^2$): fark $\le0.17$.
\item \emph{Tek demet CW yönü ve düzeltilmemiş desenler} ($10^2$--$10^4~\mu$m$^2$): fark 1--3000.
Düzeltilmemiş desenlerde eski yöntem $C$'yi 1.1--13 kat yüksek veriyor (medyan 150'ye karşı 49) ve bir
seed'de CW işareti farklı ($-0.08$'e karşı $+38$).
\end{itemize}
Bir incelik: $C$ iki yönün farkıdır. Eski yöntemin hatası iki yönde farklıdır, çünkü iki yönün kapalı
yörüngeleri farklıdır (kuadrupol tekmeleri yönle işaret değiştirir). Bu yüzden hata $C$'de sadeleşmez;
tek demet düzeltmesinde olduğu gibi bir yönün yörüngesi büyük, diğerininki küçükse hata doğrudan $C$'ye
geçer~\YOR.}


\title{Sahte EDM Ölçümünde Fırlatma Yöntemi:\\ Makaledeki Yöntem ile Donmuş Spin Yönteminin Karşılaştırılması}
\author{Hazırlayan: Claude (Anthropic yapay zekâ asistanı)}
\date{13 Kasım 2026}

\begin{document}
\sloppy
\maketitle

\begin{abstract}
\noindent Sahte EDM'yi C++ izleyicide ölçmek için parçacığın nasıl fırlatıldığını iki yöntemle
karşılaştırıyorum. \emph{Eski yöntem} (makaledeki): parçacık kapalı yörüngenin üzerinden ve fırlatma
noktasında momentum büyüklüğü tam sihirli değerde ($|p|=p_{\rm sihirli}$) fırlatılır. \emph{Yeni yöntem}:
parçacığın enerjisi, yatay spin presesyonu sıfır olacak biçimde ayarlanır ($\omega_h=0$, ``donmuş spin''),
yani deneyin kendisinin sabit tuttuğu koşulla. İki yöntem aynı kapalı yörünge aramasını ve aynı eğim
kestirimini kullanır; tek fark, hangi enerjide ölçüldüğüdür. Bulgular: (1)~eski yöntemin sonucu, aynı
fizik için fırlatma noktasına göre $-456$ ile $+2729\times$hedef arasında değişir, yeni yöntemin sonucu
\%0.2 içinde aynıdır; (2)~nedeni, sert kenarlı saptırıcı modelinde $|p|$'nin aynı kapalı yörünge
boyunca hücreden hücreye $\sim10^{-7}$ değişmesidir; (3)~eski yöntemin hatası fırlatma anındaki kapalı
yörüngenin büyüklüğüyle orantılıdır, bu yüzden makalenin iki demetli sonuçları (yörünge zaten küçük)
neredeyse etkilenmez, ama düzeltilmemiş ya da tek demetle düzeltilmiş durumlar 1.1--35 kat yanlış
(çoğunlukla abartılmış) ölçülmüştür; (4)~$\sigma^2$ ölçekleme testi bu hatayı yakalayamaz, çünkü hata da $\sigma^2$ ile ölçeklenir.
Simülasyonlar yalnız yöntemleri karşılaştırmak için seçildi; tam kapsamlı değildir.
\end{abstract}

\tableofcontents
\newpage

%======================================================================
\section{Ne ölçüyoruz ve fırlatma neden önemli?}
%======================================================================

\paragraph{Ölçülen büyüklük.} Sahte EDM, spinin düşey bileşeninin büyüme hızıdır: $f=\mathrm{d}S_y/\mathrm{d}t$.
Hedef $f\approx1$~nrad/s'dir. C++ izleyici bir protonu ve spinini halkada birkaç yüz tur izler; her turda
belirli bir noktada $S_y$ kaydedilir ve bu dizinin eğimi ölçülür.

\paragraph{Kapalı yörünge üzerinden fırlatma (iki yöntemde ortak).} Kuadrupol kaçıklıkları kapalı
yörüngeyi ideal eksenden kaydırır. Parçacık ideal eksenden fırlatılırsa kapalı yörünge etrafında
betatron salınımı yapar; bu salınımın $S_y$'ye katkısı aradığımız seküler eğimden yaklaşık bin kat
büyüktür. Bu yüzden iki yöntemde de önce 4 boyutlu kapalı yörünge ($x,x',y,y'$) aranır ve parçacık onun
üzerinden fırlatılır~\BEL.

\paragraph{Eğim kestirimi (iki yöntemde ortak).} $S_y(t)$'ye ``doğru $+$ salınımlar'' modeli fit edilir:
$S_y(t)=a+bt+\sum_k[c_k\cos(2\pi f_kt)+d_k\sin(2\pi f_kt)]$; yalnız $b$ raporlanır. Düz doğru fiti
kullanılmaz, çünkü kalan küçük salınımlar eğimi bozar.

\paragraph{Fırlatmada açık kalan tek seçim: enerji.} Kapalı yörünge konumu ve açısı belirler; ama
parçacığın \emph{enerjisini} (momentum büyüklüğünü) belirlemez. Her enerji için ayrı bir kapalı yörünge
vardır. İki yöntem tam olarak burada ayrılır. Bu önemsiz bir ayrıntı değildir: sahte EDM enerjiye çok
duyarlıdır (Bölüm~\ref{sec:delta}).

%======================================================================
\section{Simülasyonun geometrisi}
%======================================================================

Bir FODO hücresi sekiz elemandan oluşur ve C++ izleyici her turu \emph{sürüklenme (drift)} elemanı 1'den
başlatır:
\begin{kodblok}
// integrator.cpp, run_integration
// Start each revolution from the first drift element (elem=1),
// then wrap around: 1,2,3,4,5,6,7,0.
const int start_elem = 1;
...
if (elem == 0 || elem == 4) { type = 0; target_val = Phi_def; }        // ARC1, ARC2: saptırıcı
else if (elem == 1 || elem == 3 || elem == 5 || elem == 7) { type = 1; }  // sürüklenme
else if (elem == 2) { type = (current_fodo == 0) ? 4 : 2; }               // QF
else if (elem == 6) { type = 3; }                                          // QD
\end{kodblok}
Yani fırlatma noktası ARC1 saptırıcısının hemen çıkışı, QF'ten öncedir: \textbf{alansız bir bölge}.

Saptırıcıların alanı sert kenarlıdır:
\begin{kodblok}
if (element_type == 0) {
    // DEFLECTOR: cylindrical capacitor, hard-edge (no fringe fields).
    ...
    if (n == 1.0) { E_r = E0 * R0 / R; }
\end{kodblok}
Alan saptırıcının içinde $E_r=E_0R_0/R$, dışında sıfırdır; sınırda saçak alanı ve potansiyel basamağı
yoktur. Bunun sonucu Bölüm~\ref{sec:p}'de önemli olacak.

\sekil[0.95]{ek_sekil_s4_firlatma.pdf}{Üst: bir FODO hücresinin eleman sırası ve fırlatma noktası. Sol alt:
aynı kapalı yörüngenin farklı hücrelerindeki $|p|$ (Bölüm~\ref{sec:p}). Sağ alt: aynı desen, üç fırlatma
kuralı, dört farklı fırlatma noktası (Bölüm~\ref{sec:degismezlik}).}

%======================================================================
\section{Eski yöntem (makaledeki)}
%======================================================================

\subsection{Adımlar}
\begin{enumerate}[nosep]
\item Parametreler \kod{params.json}'dan okunur; sihirli momentum $p_{\rm sihirli}$ hesaplanır.
\item 4 boyutlu kapalı yörünge aranır: fırlatma noktasında $|p|=p_{\rm sihirli}$ olan bir parçacığın
betatron salınım varyansı Newton yöntemiyle en küçük yapılır.
\item Parçacık bulunan $(x,x',y,y')$ noktasından, \textbf{momentum büyüklüğü $p_{\rm sihirli}$} ve spin
boylamsal olarak fırlatılır.
\item $3\times10^{-4}$~s (yaklaşık 67 tur) izlenir; $S_y$'nin eğimi model fitiyle ölçülür.
\end{enumerate}
\textbf{Sabit tutulan:} fırlatma noktasındaki momentum büyüklüğü.

\subsection{Kod}
Makalenin Şekil~4'ünü üreten \kod{kmod\_drivers/twobeam\_oc.py} her ölçüm için şunu çağırır:
\begin{kodblok}
# kmod_drivers/twobeam_oc.py
def _C(dx, dy, dG=None, tilt=None):
    fcw, _ = fast_measure(dx, dy, tilt=tilt, direction=+1, signed=True, n_turns=28, n_iter=4, dG=dG)
    fccw, _ = fast_measure(dx, dy, tilt=tilt, direction=-1, signed=True, n_turns=28, n_iter=4, dG=dG)
    return fcw / TARGET, fccw / TARGET
\end{kodblok}
\kod{fast\_measure}'ın çekirdeği (\kod{kmod\_drivers/fast\_est.py}):
\begin{kodblok}
fields, y0, beta0, R0, p_mag, direction = setup_fields(cfg)        # p_mag = p_sihirli
T_rev = _T_rev(fields, beta0, R0)
v_co, resid = find_co_4d(fields, p_mag, direction, dx, dy, tilt, T_rev,
                         n_turns=n_turns, n_iter=n_iter, dG=dG)   # adım 2
launch = _make_state(v_co, p_mag, direction, [0.0, 0.0, direction])  # adım 3
fields.poincare_quad_index = 0.0
_, poin, pt = integrate_particle(launch, 0.0, t2, DT, fields=fields, ...)   # adım 4
slope = float(measure_dSy_dt_model(poin[:, 7], pt))
\end{kodblok}
Fırlatma durumunu kuran fonksiyon (\kod{berry\_data/false\_edm\_mode\_scan.py}); momentum büyüklüğü
doğrudan \kod{p\_mag}'dir:
\begin{kodblok}
def _make_state(v, p_mag, direction, spin):
    """(x0, y0, x', y') faz-uzayı noktasından 9-bileşenli yerel başlangıç durumu."""
    return [v[0], v[1], 0.0,
            p_mag * direction * v[2], p_mag * direction * v[3], p_mag * direction,
            spin[0], spin[1], spin[2]]
\end{kodblok}
Kapalı yörünge araması (\kod{berry\_data/false\_edm\_4d.py}); aranan büyüklük her turda aynı noktadaki
konumun varyansıdır ve aramada da aynı \kod{\_make\_state} kullanılır:
\begin{kodblok}
def find_co_4d(fields, p_mag, direction, dx, dy, tilt, T_rev, n_turns=28, n_iter=2, dG=None):
    def var4(v):
        st = _make_state(v, p_mag, direction, spin)
        _, poin, _ = integrate_particle(st, 0.0, t_probe, DT, fields=fields, ...)
        return float(np.var(poin[:, 0]) + np.var(poin[:, 1]))
    # 4x4 Hessian ile Newton adımları (n_iter kez)
\end{kodblok}

%======================================================================
\section{Yeni yöntem (donmuş spin)}
%======================================================================

\subsection{Fikir}
Gerçek deneyde demetin enerjisi ne momentumla ne toplam enerjiyle sabitlenir; polarimetre spinin yatay
yönünü (fazını) ölçer ve bir geri besleme, yatay presesyonu sıfır tutacak biçimde enerjiyi ayarlar
(``donmuş spin''). JEDI işbirliği COSY'de presesyon fazını geri beslemeyle kilitlemeyi göstermiştir. Yeni
yöntem simülasyonda aynı koşulu kurar: sahte EDM, $\omega_h=0$ olan enerjide ölçülür.

\subsection{Adımlar}
\begin{enumerate}[nosep]
\item Bir başlangıç kuralıyla ($\delta=0$) kapalı yörünge aranır ve izlenir: $f_0$ ve yatay presesyon hızı
$\omega_0=\mathrm{d}/\mathrm{d}t\,\mathrm{atan2}(S_x,S_z)$ ölçülür.
\item Enerji düzeltilir: $\delta_1=-\omega_0/(\mathrm{d}\omega_h/\mathrm{d}\delta)$. Eğim
$\mathrm{d}\omega_h/\mathrm{d}\delta=-5.238\times10^6$~rad/s kaçıklıktan bağımsızdır (Bölüm~\ref{sec:delta}).
Bu $\delta_1$ ile kapalı yörünge \emph{yeniden} aranır ve izlenir: $f_1$, $\omega_1$.
\item $f$ ve $\omega_h$, $\delta$'ya doğrusal olduğu için $\omega_h=0$'daki değer ara değerlemeyle alınır:
$f^*=f_0+(f_1-f_0)\,(0-\omega_0)/(\omega_1-\omega_0)$.
\end{enumerate}
\textbf{Sabit tutulan:} $\omega_h=0$. Başlangıç kuralı yalnız bir başlangıç tahminidir; sonucu etkilemez
(Bölüm~\ref{sec:baslangic}).

\subsection{Kod}
Tamamı \kod{revizyon/sahte\_edm\_donmus.py}'dedir. Başlangıç kuralı (enerji korunumlu tahmin + $\delta$):
\begin{kodblok}
def launcher(delta):
    def make_state(v, p_mag, direction, spin):
        x = v[0]
        W_ref = np.sqrt((p_mag * C_LIGHT) ** 2 + (M_P * C_LIGHT ** 2) ** 2)
        Wk = W_ref - Q_E * (-E0_MAGIC * R0 * np.log((R0 + x) / R0))
        p = np.sqrt(Wk ** 2 - (M_P * C_LIGHT ** 2) ** 2) / C_LIGHT * (1.0 + delta)
        return [v[0], v[1], 0.0, p * direction * v[2], p * direction * v[3], p * direction,
                spin[0], spin[1], spin[2]]
    return make_state
\end{kodblok}
Tek bir $\delta$ için ölçüm: eski yöntemle aynı \kod{find\_co\_4d}, \kod{integrate\_particle} ve
\kod{measure\_dSy\_dt\_model}; tek ek, yatay presesyon hızının da ölçülmesi:
\begin{kodblok}
def _one(dx, dy, direction, delta, tilt, dG, n_turns, n_iter, t2):
    ms = launcher(delta)
    false_edm_4d._make_state = ms; false_edm_mode_scan._make_state = ms   # aramada da aynı kural
    fields, y0, beta0, R0_, p_mag, d = setup_fields(cfg)
    v_co, resid = find_co_4d(fields, p_mag, d, dx, dy, tilt, _T_rev(...), n_turns=n_turns, n_iter=n_iter, dG=dG)
    launch = ms(v_co, p_mag, d, [0.0, 0.0, d])
    _, poin, pt = integrate_particle(launch, 0.0, t2, dt, fields=fields, ...)
    f = measure_dSy_dt_model(poin[:, 7], t)                               # eskisiyle aynı
    ang = np.unwrap(np.arctan2(poin[:, 6], poin[:, 8] * d))              # yatay spin açısı
    return f, np.polyfit(t, ang, 1)[0], resid                            # f, omega_h
\end{kodblok}
Donmuş spin noktası:
\begin{kodblok}
def measure_frozen(dx, dy, direction=-1, tilt=None, dG=None, n_turns=28, n_iter=4, t2=3e-4):
    f0, w0, r0 = _one(dx, dy, direction, 0.0, tilt, dG, n_turns, n_iter, t2)
    d1 = -w0 / S_W_DELTA                                    # S_W_DELTA = -5.238e6 rad/s
    f1, w1, r1 = _one(dx, dy, direction, d1, tilt, dG, n_turns, n_iter, t2)
    fs = f0 + (f1 - f0) * (0.0 - w0) / (w1 - w0)           # omega_h = 0'da f
    return dict(f=fs, f0=f0, f1=f1, w0=w0, w1=w1, delta_star=d1 * (0.0 - w0) / (w1 - w0), ...)
\end{kodblok}

\subsection{Yan yana}
\begin{center}\small
\begin{tabularx}{\textwidth}{>{\raggedright\arraybackslash}p{3.6cm}>{\raggedright\arraybackslash}X>{\raggedright\arraybackslash}X}
\toprule
& eski (makale) & yeni (donmuş spin) \\
\midrule
kapalı yörünge araması & \kod{find\_co\_4d} & aynı \\
eğim kestirimi & \kod{measure\_dSy\_dt\_model} & aynı \\
sabit tutulan & fırlatma noktasında $|p|=p_{\rm sihirli}$ & $\omega_h=0$ (deneydeki koşul) \\
izleme sayısı & 1 (+ kapalı yörünge araması) & 2 (+ iki kapalı yörünge araması) \\
maliyet & 1 birim & $\approx2$ birim \\
ek çıktı & -- & $\omega_0$, $\omega_1$, $\delta^*$ (donmuş noktanın enerjisi) \\
\bottomrule
\end{tabularx}
\end{center}

%======================================================================
\section{Karşılaştırma testleri}
%======================================================================
Bütün testler $\sigma=10~\mu$m (rng 3000, seed~0) deseninde ya da makalenin Şekil~4 kurulumunda (5 seed)
yapıldı. Ayar: kapalı yörünge araması 28 tur, 4 iterasyon; izleme $3\times10^{-4}$~s.

\subsection{Test 1: Fırlatma noktasından bağımsızlık}\label{sec:degismezlik}
Halka simetrik olduğu için, aynı kaçıklık desenini halkada $k$ hücre kaydırmak yalnızca fırlatma
noktasını değiştirir; fizik aynıdır. Doğru bir yöntem dört fırlatma noktasında aynı sonucu vermelidir
(\kod{tests/teshis\_enerji\_ABC.py})~\BEL:
\begin{center}\small
\begin{tabular}{lcccc}
\toprule
fırlatma hücresi & 0 & 6 & 12 & 18 \\
\midrule
eski: $|p|=p_{\rm sihirli}$ & 1088 & 2729 & $-456$ & 95 \\
ara deneme: ``toplam enerji tasarımda'' & $-376$ & 162 & $-722$ & 567 \\
\textbf{yeni: donmuş spin} & \textbf{82.04} & \textbf{81.90} & \textbf{81.89} & \textbf{81.94} \\
\bottomrule
\end{tabular}
\end{center}
(Hedefin katı.) Eski yöntem aynı fizik için $-456$ ile $+2729$ arasında değişiyor; yeni yöntem \%0.2 içinde
aynı. Şekil~1'in sağ alt paneli bunu gösteriyor.

\subsection{Test 2: Neden? Aynı yörüngede $|p|$ hücreden hücreye değişiyor}\label{sec:p}
Eski yöntemle hücre 0'dan fırlatılan parçacığın kapalı yörüngesi boyunca, hücre 0, 6, 12, 18'deki $|p|$
doğrudan ölçüldü (\kod{tests/teshis\_p\_azimut.py}; ölçüm ARC2 girişinde, yani fırlatma noktasıyla
arasında yalnız sürüklenme ve kuadrupol olan yerde)~\BEL:
\begin{center}\small
\begin{tabular}{lcccc}
\toprule
hücre & 0 & 6 & 12 & 18 \\
\midrule
$\sigma=10~\mu$m: $|p|/p_{\rm sihirli}-1$ & $+2\times10^{-12}$ & $-5.48\times10^{-7}$ & $-2.82\times10^{-7}$ & $-2.33\times10^{-7}$ \\
kaçıklık yok & $\sim10^{-12}$ & $\sim10^{-12}$ & $\sim10^{-12}$ & $\sim10^{-12}$ \\
\bottomrule
\end{tabular}
\end{center}
Mekanizma~\YOR: sert kenarlı saptırıcının içinde elektrik alanı parçacığın radyal hareketine göre iş yapar.
Gerçek bir saptırıcıda bu iş, saçak alanındaki potansiyel basamağında geri alınır ve toplam enerji korunur.
Sert kenarlı modelde basamak yoktur; bu yüzden kapalı yörüngenin saptırıcı içindeki radyal değişimine bağlı
olarak $|p|$ bir hücreden diğerine değişir. Kaçıklık yokken kapalı yörünge ideal eksendir ve değişim
sıfırdır.

Sonuç: eski yöntem ``hücre $k$'da $|p|=p_{\rm sihirli}$'' derken, her $k$ için $\sim10^{-7}$ farklı bir
yörünge enerjisi seçer.

Çapraz kontrol~\BEL: Test~1'deki dört donmuş spin çözümünden, donmuş spin yörüngesinin her hücredeki
$|p|$'si hesaplandı; hücre 0'a göre farklar $-5.482$, $-2.816$, $-2.333\times10^{-7}$. Doğrudan ölçümle
dört basamak aynı. Yani dört farklı noktadan bulunan donmuş spin çözümleri \emph{tek ve aynı yörüngedir}.

\subsection{Test 3: Sahte EDM enerjiye ne kadar duyarlı?}\label{sec:delta}
Aynı desende enerji $\delta=\pm10^{-7},\pm10^{-6}$ kaydırılarak izlendi~\BEL:
\sekil[0.9]{ek_sekil_s3_delta.pdf}{Enerji taraması. Sol: yatay presesyon; sağ: sahte EDM. İkisi de $\delta$'ya
doğrusal; elmas, donmuş spin noktası.}
$\delta=10^{-7}$ başına sahte EDM $\approx290\times$hedef değişiyor. Test~2'deki $|p|$ farkları
($2$--$5\times10^{-7}$) bu yüzden yüzlerce hedeflik farklar yaratır. Ayrıca~\BEL\ (donmuş spin aracının 9
koşumundan): $\mathrm{d}\omega_h/\mathrm{d}\delta=-5.238\times10^6$~rad/s kaçıklıktan bağımsız (halkanın
özelliği), $\mathrm{d}f/\mathrm{d}\delta$ ise kaçıklıkla doğrusal ölçekleniyor.

\subsection{Test 4: Yeni yöntem başlangıç kuralından bağımsız mı?}\label{sec:baslangic}
Yeni yöntemin bir başlangıç kuralına ihtiyacı var ($\delta=0$ noktası). Varsayılan kural ``enerji
korunumlu'' tahmindir ve bu, fırlatma noktası alansız bölgede olduğu için fiziksel bir anlam taşımaz. Bu
yüzden aynı prosedür, başlangıç kuralı olarak \emph{eski yöntemin kuralıyla} ($|p|=p_{\rm sihirli}(1+\delta)$)
dört fırlatma hücresinde tekrar çalıştırıldı (\kod{tests/teshis\_baslangic\_kurali.py})~\BEL:
\BaslangicTablo
\sekil[0.9]{firlatma_fD_baslangic.pdf}{Sol: eski kuraldan başlayınca başlangıç değeri fırlatma noktasına göre
çok farklı, ama donmuş spin sonucu aynı. Sağ: iki başlangıç kuralının verdiği donmuş spin sonuçları (dar eksen).}
\BaslangicYorum

\subsection{Test 5: $\sigma^2$ ölçekleme testi bu farkı neden göremez?}
Makalede ve önceki çalışmalarda eski yöntem, $\sigma=2.5$, 5, 10~$\mu$m'de $p\approx2.00$ üssüyle
doğrulanmıştı. Yeni yöntem de $p=2.0002$ veriyor. Ama iki yöntemin değerleri farklı (10~$\mu$m'de eski
1089, 58, 1587; yeni 82, $-71$, 544).

Nedeni~\BEL: eski yöntemin hatası $=(\mathrm{d}f/\mathrm{d}\delta)\times\delta_{\rm fırlatma}$. İki çarpan da
kaçıklıkla doğrusal: $\mathrm{d}f/\mathrm{d}\delta\propto\sigma$ (Test~3), $\delta_{\rm fırlatma}\propto
x_{\rm CO}\propto\sigma$ (Test~2). Çarpım $\propto\sigma^2$, yani hata doğru değerle \emph{aynı üsle}
ölçeklenir.
\sekil[0.9]{ek_sekil_s5_sigma2.pdf}{Sol: sabit başlangıç kuralı ve donmuş spin; ikisi de $\sigma^2$. Sağ:
oranları her seed için $\sigma$'dan bağımsız ama 1'den farklı ($-4.6$, $2.4$, $0.84$).}
Ders: $\sigma^2$ testi mekanizmanın türünü (ikinci mertebe geometrik faz) doğrular; mutlak değeri
doğrulamaz. Mutlak değer için Test~1 gibi bir değişmezlik testi gerekir.

\subsection{Test 6: Makalenin Şekil 4'ü, aynı desenlerle}\label{sec:fig4}
Eski değerler makalenin kendi çıktısıdır (\kod{kmod\_drivers/twobeam\_oc\_14nm.json}); yeni değerler aynı
desen, aynı gürültü, aynı düzeltmeyle donmuş spin aracından (\kod{donmus\_fig4.jsonl})~\BEL:
\sekil[0.95]{firlatma_fB_sekil4.pdf}{Şekil 4 kurulumu (10~$\mu$m korelasyonsuz kaçıklık, 14~nm gürültü,
\kod{rcond}~$=0.01$). $C$, seed başına.}
\begin{center}\small
\begin{tabular}{lcccccc}
\toprule
& \multicolumn{3}{c}{eski} & \multicolumn{3}{c}{yeni} \\
\cmidrule(lr){2-4}\cmidrule(lr){5-7}
durum & $f_{\rm CW}$ medyan & $f_{\rm CCW}$ medyan & $C$ medyan & $f_{\rm CW}$ & $f_{\rm CCW}$ & $C$ \\
\midrule
düzeltme yok & \HamEskiCW & \HamEskiCCW & \HamEskiC & 38 & 78 & 49 \\
tek demet & $-25.6$ & $-1.55$ & 47 & $-10.1$ & $-1.56$ & 12 \\
iki demet & 0.006 & 0.001 & 0.0046 & 0.000 & 0.002 & 0.0015 \\
\bottomrule
\end{tabular}
\end{center}
Tablo yanıltıcı olmasın diye her yönü ayrı inceledim. Tek demet düzeltmesinde yalnız CCW yörüngesi
düzeltilir; CW yörüngesi büyük kalır. Eski ve yeni yöntem CCW'de neredeyse aynı (fark $\le0.17$), CW'de çok
farklı (fark 1--100)~\BEL. Bu, hatanın kapalı yörüngenin büyüklüğüne bağlı olduğunu düşündürür; doğrudan
sınadım:
\sekil[0.8]{firlatma_fA_fark_yorunge.pdf}{Eski ve yeni yöntem arasındaki fark, fırlatma anındaki kapalı
yörüngenin büyüklüğüne (yatay ve düşey rms'in çarpımı, $R$ matrisinden) karşı. Her nokta bir seed ve bir
yön.}
\FarkYorum

\subsection{Test 7: Önceki raporlardaki ana sayılar}
Revizyon kontrolleri sırasında eski yöntemle elde edilip yeni yöntemle tekrarlanan sayılar~\BEL:
\sekil[0.95]{firlatma_fC_ozet.pdf}{Her satır bir sonuç; turuncu eski, mavi yeni yöntem. Elmaslar medyan,
noktalar seed'ler. Eski değerler 3--5 seed, yeniler 5--10 seed.}
\begin{center}\small
\begin{tabular}{lccl}
\toprule
sonuç & eski medyan & yeni medyan & yorum \\
\midrule
Şekil 4, iki demet & 0.0046 & 0.0015 & ikisi de tabanda; değişmedi \\
Şekil 4, tek demet & 47 & 12 & 4 kat abartılmıştı \\
Adım 1: $\rho=1$, makale kestiricisi & 2.5 & 0.19 & hüküm değişti: hedefin altında \\
T2: rijit BPM, iki demet, $\sigma_\delta=1~\mu$m & 0.27 & 0.010 & 27 kat abartılmıştı \\
Adım 6: makale modeli, 96. saat & 15 & 1.3 & ``çalışmaz'' $\to$ ``marjsız'' \\
T1: yalnız kuadrupol, \%0.2 $\beta$-beat & 1.3 & 0.75 & aynı yönde, 2 kat \\
\bottomrule
\end{tabular}
\end{center}
Desen tutarlı~\YOR: kalan kapalı yörünge küçükse (iki demet, trim) iki yöntem uyuşuyor; büyükse (tek demet,
yanlış BPM modeli, düzeltilmemiş) eski yöntem sahte EDM'yi ya abartıyor ya da yanlış işaretle veriyor.
Hükmü değişen iki sonuç var: Adım~1'de makale kestiricisi $\rho=1$'de hedefin altına iniyor; Adım~6'da
makale modeli ``çalışmaz'' değil ``marjsız''.

%======================================================================
\section{Hangisi doğru, ve ne kadar emin olabiliriz?}
%======================================================================

\paragraph{Yeni yöntemden yana kanıtlar~\YOR.}
\begin{enumerate}[nosep]
\item \emph{Deneyle aynı çalışma noktası:} deney $\omega_h=0$'ı sabit tutar; yeni yöntem de.
\item \emph{Değişmezlik:} fiziksel olarak aynı durum (farklı fırlatma noktası) aynı sonucu veriyor (Test~1);
eski yöntem vermiyor.
\item \emph{Tek yörünge:} dört noktadan bulunan çözümler aynı yörünge (Test~2 çapraz kontrolü).
\item \emph{Başlangıçtan bağımsızlık:} başlangıç kuralı değişse de sonuç aynı (Test~4).
\end{enumerate}

\paragraph{Eski yöntem ``yanlış fizik'' mi?} Hayır~\YOR. Eski yöntem doğru bir izleme yapıyor, ama kontrol
edilmeyen bir enerjide. Ölçtüğü değer, o enerjideki sahte EDM'dir; o enerji ise fırlatma noktasına ve kapalı
yörüngeye bağlıdır. Sorun ölçümde değil, çalışma noktasının tanımındadır.

\paragraph{Saçak alanlı (gerçekçi) bir modelde eski yöntem doğru olur muydu?~\HIP} Kısmen. Saçak alanı olan
bir saptırıcıda toplam enerji korunur ve alansız bölgelerin hepsinde $|p|$ aynıdır; Test~1'deki fırlatma
noktası bağımlılığı büyük ölçüde kaybolurdu. Ama kaçıklık, kapalı yörüngenin uzunluğunu ve saptırıcılardaki
ortalama yarıçapını değiştirir; bu da $\omega_h=0$ olan enerjiyi tasarım enerjisinden kaydırır. Eski yöntem
tasarım enerjisinde ölçeceği için yine donmuş noktanın dışında kalırdı. Bu kaymanın büyüklüğünü bu modelde
ölçemiyorum (sert kenar onu fırlatma noktası etkisiyle karıştırıyor). Her iki durumda da deneyle eşleşen
tanım $\omega_h=0$'dır.

\paragraph{Yeni yöntemin sınırları.}
\begin{itemize}[nosep]
\item Sert kenarlı saptırıcı modeli değişmedi; yeni yöntem yalnız bu modelde çalışma noktasını doğru seçer.
Saçak alanının kendisinin sahte EDM'ye etkisi test edilmedi.
\item $f(\delta)$ ve $\omega_h(\delta)$'nın doğrusallığı tek desende ($\delta=\pm10^{-6}$) sınandı. İkinci nokta
donmuş noktadan uzak düşerse iki noktalı ara değerleme \%3'e kadar sapabilir (Test~4); araç artık bu durumda
$\delta^*$'da üçüncü bir ölçüm yapıyor. Önceki 666 koşumda ikinci nokta zaten donmuş noktadaydı~\BEL.
\item Geri beslemenin dinamiği ve demetin momentum dağılımı ($\sim10^{-4}$) simüle edilmedi; tek parçacık
ortalama $\delta$'yı temsil eder. $f$ ve $\omega_h$ $\delta$'ya doğrusal olduğu için ortalama alma bunu
değiştirmez~\YOR.
\item Maliyet yaklaşık iki katı.
\item Eski yöntemin donmuş nokta dışındaki $\mathrm{d}f/\mathrm{d}\delta$ eğimi fırlatma noktasına göre
farklı çıkıyor; bunu açıklamadım. Bu, eski değerleri tek bir eğimle ``düzeltmeyi'' imkânsız kılıyor:
eski sayılar yeniden ölçülmelidir, dönüştürülemez.
\end{itemize}

%======================================================================
\section{Makale için sonuç}
%======================================================================
\begin{enumerate}[nosep]
\item Makalenin ana sayısı, iki demetli düzeltmeden sonraki sahte EDM, iki yöntemde de sayısal tabanda
(0.0046 ve 0.0015). Ana iddia etkilenmiyor~\BEL.
\item Karşılaştırma için verilen sayılar (düzeltme yok, tek demet) eski yöntemle 1.1--35 kat yanlış~\BEL.
Makalede ``iki demet tek demetten $\sim10^4$ kat iyi'' türünden bir oran varsa, yeni yöntemle yaklaşık
$8\times10^3$ olur (12 / 0.0015); oranın büyüklük mertebesi değişmiyor, ama tek demetin kötülüğünün büyük
kısmı \kod{rcond}~$=0.01$ seçiminden geliyor (ek rapor, soru 12).
\item Makalede ölçüm prosedürü tanımlanırken ``parçacık kapalı yörüngeden, $\omega_h=0$ olan enerjide
fırlatılır'' yazılmalı ve kısa bir değişmezlik testi (Test~1) verilmelidir. Bu, hakemin olası ``hangi
enerjide?'' sorusunu önceden cevaplar.
\item Sert kenar modelinin bir sınırlama olarak yazılması önerilir.
\end{enumerate}

%======================================================================
\appendix
\section{Yeniden üretim}
%======================================================================
Dal: \kod{claude/brave-allen-92wn61}.
\begin{center}\small
\begin{tabularx}{\textwidth}{>{\raggedright\arraybackslash}p{8.6cm}X}
\toprule
dosya & işlev \\
\midrule
\kod{kmod\_drivers/fast\_est.py},\newline \kod{berry\_data/false\_edm\_4d.py},\newline \kod{berry\_data/false\_edm\_mode\_scan.py} & eski yöntem (değiştirilmedi) \\
\kod{revizyon/sahte\_edm\_donmus.py} & yeni yöntem \\
\kod{revizyon/tests/teshis\_enerji\_ABC.py} & Test 1 ve 3 \\
\kod{revizyon/tests/teshis\_p\_azimut.py} & Test 2 \\
\kod{revizyon/tests/teshis\_baslangic\_kurali.py},\newline \kod{revizyon/tests/teshis\_baslangic\_dogrulama.py} & Test 4 \\
\kod{revizyon/donmus\_fig4.py},\newline \kod{revizyon/firlatma\_eski\_ham.py} & Test 6 \\
\kod{revizyon/firlatma\_sekiller.py},\newline \kod{revizyon/ek\_sekiller.py} & şekiller \\
\bottomrule
\end{tabularx}
\end{center}
\end{document}
